\documentclass[11pt]{article}
\usepackage[T1]{fontenc}
\usepackage{lmodern}
\usepackage{amsmath,amssymb,amsthm,booktabs,longtable}
\usepackage[margin=1in]{geometry}
\usepackage[colorlinks=true,urlcolor=blue,citecolor=blue,linkcolor=blue]{hyperref}
\urlstyle{same}
\newtheorem{theorem}{Theorem}
\newtheorem{lemma}{Lemma}
\newcommand{\Z}{\mathbb Z}
\newcommand{\R}{\mathbb R}
\newcommand{\C}{\mathbb C}
\DeclareMathOperator{\Spin}{Spin}
\title{A type-$D_4$ counterexample to the\\canonical-potential exponent converse}
\author{Alec Kriebel\\
\small Independent researcher\\
\small\href{https://orcid.org/0009-0001-9320-500X}{ORCID: 0009-0001-9320-500X}}
\date{17 October 2026 \quad Version 1.0}
\begin{document}
\maketitle
\begin{abstract}
The exponent-bound converse stated for canonical GHKK potentials in Schumann's
2021 Oberwolfach report is false. For an explicit reduced word in type $D_4$,
the canonical potential has $22$ combined monomials, all of coefficient one,
and all $22$ inequalities define distinct facets, although two monomials have
exponent $-2$ in one canonical coordinate. The reduced word and its complete
string-coordinate support occur in earlier public computations; we extract
the consequence for the historical formulation and supply an independent
canonical calculation and exact integer certificates. The chamber-ansatz
pullback has all exponents in $\{-1,0,1\}$, so the later string-coordinate
converse is not refuted. No first-priority or current-openness claim is made.
\end{abstract}
\noindent\textbf{Status.} Extensively AI-assisted, unrefereed research note.
There has been no external human peer review or proof-assistant certification.

\section{The historical assertion and the result}
Schumann's report \cite[printed p.~1125]{OWR} formulates a converse for the
canonical potential $W$ restricted to the dual cluster torus designated by
a reduced word $\mathbf i$ for the longest Weyl group element. In that
formulation, the monomial inequalities are irredundant precisely when every
coordinate exponent of $W$ has absolute value at most one. The correspondence
with the string cone is unimodular \cite{GKS,KS}. An invertible monomial
change preserves irredundancy, but need not preserve that exponent bound.

We use the reduced-word seed and $\mathcal X$-mutation conventions of
\cite[Eq.~(3), Defs.~7--9]{KS}, including Dynkin adjacency in the inclined-arrow
rule as specified in \cite[\S4.2.2]{GKS}. These are the designated canonical
coordinates, not an arbitrary change chosen to create a large exponent.

\begin{theorem}\label{thm:main}
Let $G=\Spin_8(\C)$, with Dynkin edges $1\!-\!2$, $2\!-\!3$, $2\!-\!4$, and
\[
 \mathbf i=(1,2,3,2,4,2,1,2,3,2,4,2).
\]
In its canonical reduced-word $\mathcal X$-coordinates, $W$ has $22$ distinct
monomials, each with coefficient one. Their $22$ tropical inequalities define
distinct facets of a full-dimensional cone in $\R^{12}$. Two monomials
have $X_6$-exponent $-2$. Consequently the necessary direction of the
canonical-$W$ converse in \cite{OWR} fails.
\end{theorem}

The word and its complete chamber-pullback support already occur in the
earlier computational data \cite[Case~267]{DATA}; see Section~\ref{sec:priority}.
We do not claim first computation of these objects. The theorem is confined
to the historical canonical-coordinate assertion. It does not refute the
later $\varsigma$-coordinate converse or coefficient conjecture in
\cite{KS}.

\section{The canonical potential}
Put $y_k=X_k^{-1}$. The frozen positions are $7,9,11,12$. The complete arrow
set, with all multiplicities one, is
\[
\begin{gathered}
 1\to7,\quad2\to4,\quad3\to2,\quad3\to9,\quad4\to6,\quad5\to4,\quad5\to11,\quad6\to1,\\
 6\to8,\quad7\to6,\quad8\to3,\quad8\to10,\quad9\to8,\quad10\to5,\quad10\to12,\quad11\to10.
\end{gathered}
\]
For completeness, let $k^+$ be the next occurrence of $i_k$, or $13$.
Horizontal arrows are $k\to k^+$ when $k^+\le12$. There is an inclined
arrow $\ell\to k$ for $k<\ell<k^+<\ell^+$ when $i_k,i_\ell$ are Dynkin
neighbors. Arrows between frozen vertices are omitted.

The inversion roots, in the ordered simple-root basis, are
\[
\begin{gathered}
(1,0,0,0),(1,1,0,0),(1,1,1,0),(0,0,1,0),\\
(1,1,1,1),(1,1,0,1),(1,2,1,1),(0,1,1,0),\\
(0,1,1,1),(0,0,0,1),(0,1,0,1),(0,1,0,0).
\end{gathered}
\]
These are all twelve positive $D_4$ roots, and the product of the simple
reflections is $-I_4$, so the word is reduced for $w_0$.

Set $b_{ab}=\#(a\to b)-\#(b\to a)$. Mutation at a mutable position $k$ gives
\begin{equation}\label{eq:mutation}
 y'_k=y_k^{-1},\qquad
 y'_a=\begin{cases}
 y_a(1+y_k)^{b_{ak}}&b_{ak}\ge0,\ a\ne k,\\
 y_a y_k^{-b_{ak}}(1+y_k)^{b_{ak}}&b_{ak}<0,\ a\ne k.
 \end{cases}
\end{equation}
The exchange matrix changes by sign reversal in row/column $k$ and otherwise
\[
 b'_{ab}=b_{ab}+[b_{ak}]_+[b_{kb}]_+-[-b_{ak}]_+[-b_{kb}]_+,
\]
with frozen--frozen entries erased. All right-hand sides use the old matrix.
An optimized frozen position has only incoming mutable arrows; its canonical
summand is the inverse frozen coordinate in that seed \cite[Def.~7]{KS}.

\begin{lemma}\label{lem:potential}
Writing
\[
 A=1+y_3+y_2y_3,\qquad B=y_2y_3y_4y_5,
\]
the canonical potential is
\begin{align}
 F_7&=1+y_6+y_6y_8A+
 y_6y_8y_{10}\big((1+y_5)A+B(1+y_6+y_1y_6)\big),\notag\\
 F_9&=1+y_8+y_8y_{10}(1+y_5+y_4y_5),\notag\\
 W&=y_7F_7+y_9F_9+y_{11}(1+y_{10})+y_{12}. \label{eq:potential}
\end{align}
\end{lemma}
\begin{proof}
The following fixed-label sequences, read left to right, optimize the
indicated frozen position:
\[
\begin{array}{c|l}
f&\text{mutation sequence}\\ \hline
7&(2,6,1,3,4,8,5,10)\\
9&(4,8,5,10)\\
11&(10)\\
12&()
\end{array}
\]
Every entry is mutable. In each case the final frozen row has its only
nonzero entry $b_{f,10}=-1$, as follows by the displayed matrix rule.
Thus the final seed is optimized, and the canonical summand is $y_f$ there.

Pull back that monomial, in reverse order, by \eqref{eq:mutation}.
For $f=7$, after the reverse steps $10,5,8,4$ the quotient by $y_7$ is
\[
1+y_8+y_8y_{10}+y_5y_8y_{10}+y_4y_5y_8y_{10}.
\]
The next step, at $3$, gives
\[
1+y_8+y_3y_8+y_8y_{10}
 (1+y_5+y_3+y_3y_5+y_3y_4y_5).
\]
The step at $1$ adds $y_1y_3y_4y_5y_8y_{10}$. After the step at $6$ it is
\begin{align*}
1+y_6+y_6y_8(1+y_3)
&+y_6y_8y_{10}(1+y_5+y_3+y_3y_5+y_3y_4y_5)\\
&+(1+y_1)y_3y_4y_5y_6^2y_8y_{10}.
\end{align*}
The final step at $2$ gives $F_7$. The four-step $f=9$ pullback gives $F_9$;
the last two summands are immediate. These exact rational identities
identify the canonical normalized potential, rather than merely a
polynomial defining the same cone.
\end{proof}

The summands in \eqref{eq:potential} contain $14,5,2,1$ distinct terms.
Every term has exactly one frozen marker, and the coefficients are all one.
In particular
\[
 y_2y_3y_4y_5y_6^2y_7y_8y_{10},
 \qquad y_1y_2y_3y_4y_5y_6^2y_7y_8y_{10}
\]
occur. They have $X_6$-exponent $-2$.

\section{An exact certificate for every facet}
Write the $22$ terms as $y^{a^{(j)}}$. In $X$-cocharacter coordinates $t$ the
inequality is $-a^{(j)}\cdot t\ge0$. Put $z=-t$; the cone becomes
\[
 C=\{z\in\R^{12}:a^{(j)}\cdot z\ge0,\quad1\le j\le22\}.
\]
The following table gives its complete exponent support and integer vectors
$z^{(j)}$. Each exponent row is encoded as twelve consecutive single digits
in position order $1,\ldots,12$; no entry has more than one digit.

\begingroup\footnotesize
\setlength{\tabcolsep}{5pt}
\begin{longtable}{rll}
\toprule
$j$&$a^{(j)}$&$z^{(j)}$\\ \midrule
\endfirsthead
\toprule $j$&$a^{(j)}$&$z^{(j)}$\\ \midrule
\endhead
1 & \texttt{000000000001} & $(0,0,0,0,0,0,0,0,0,0,0,-1)$ \\
2 & \texttt{000000000010} & $(0,0,0,0,-1,0,0,0,0,1,-1,0)$ \\
3 & \texttt{000000000110} & $(0,0,0,0,0,0,0,1,0,-1,0,0)$ \\
4 & \texttt{000000001000} & $(0,0,-1,0,0,0,0,1,-1,0,0,0)$ \\
5 & \texttt{000000011000} & $(-2,0,0,0,0,1,0,-1,0,1,0,0)$ \\
6 & \texttt{000000011100} & $(-2,0,0,0,1,1,0,0,0,-1,1,0)$ \\
7 & \texttt{000000100000} & $(-1,0,0,0,0,1,-1,0,0,0,0,0)$ \\
8 & \texttt{000001100000} & $(0,2,0,-1,0,-1,0,1,0,0,0,0)$ \\
9 & \texttt{000001110000} & $(0,0,1,-1,0,0,0,-1,1,1,0,0)$ \\
10 & \texttt{000001110100} & $(0,0,1,-1,1,0,0,0,1,-1,1,0)$ \\
11 & \texttt{000010011100} & $(-2,0,0,1,-1,1,0,0,0,0,0,0)$ \\
12 & \texttt{000011110100} & $(0,0,1,0,-1,0,0,0,1,0,0,0)$ \\
13 & \texttt{000110011100} & $(0,1,0,-1,0,0,0,0,0,0,0,0)$ \\
14 & \texttt{001001110000} & $(0,1,-1,-1,0,0,0,0,0,1,0,0)$ \\
15 & \texttt{001001110100} & $(-1,1,-1,-1,1,1,0,0,1,-1,1,0)$ \\
16 & \texttt{001011110100} & $(-1,1,-1,0,-1,1,0,0,1,0,0,0)$ \\
17 & \texttt{011001110000} & $(0,-1,0,0,0,0,0,0,0,1,0,0)$ \\
18 & \texttt{011001110100} & $(-1,-1,0,0,1,1,0,0,1,-1,1,0)$ \\
19 & \texttt{011011110100} & $(-1,-1,0,1,-1,1,0,0,1,0,0,0)$ \\
20 & \texttt{011111110100} & $(0,0,0,-1,-1,1,0,0,2,0,0,0)$ \\
21 & \texttt{011112110100} & $(1,0,0,0,0,-1,1,0,0,0,0,0)$ \\
22 & \texttt{111112110100} & $(-1,0,0,0,0,0,0,0,0,0,0,0)$ \\
\\[-9pt]\bottomrule
\end{longtable}
\endgroup
Direct integer multiplication of the displayed rows and columns gives
\begin{equation}\label{eq:cert}
 a^{(j)}\cdot z^{(j)}=-1,\qquad
 a^{(r)}\cdot z^{(j)}\ge0\quad(r\ne j).
\end{equation}
This is a complete finite certificate: no optimization tolerance or
unprovided historical computation is a premise. Omitting the $j$th
inequality admits $z^{(j)}$, proving irredundancy.

There is also a strict facet certificate. Set
\[
 u=(0,0,0,0,0,0,1,0,1,0,1,1).
\]
Every exponent row has $a^{(j)}\cdot u=1$, so $C$ has a strict interior
point and is full-dimensional. For $p_j=z^{(j)}+u$,
\[
 a^{(j)}\cdot p_j=0,\qquad a^{(r)}\cdot p_j\ge1\quad(r\ne j).
\]
A neighborhood of $p_j$ in the hyperplane $(a^{(j)})^\perp$ satisfies all the
other inequalities strictly. Thus the selected face has dimension $11$.
The certificates separate every pair of selected inequalities, so the
$22$ facets are distinct. Together with Lemma~\ref{lem:potential}, this
proves Theorem~\ref{thm:main}.

\section{The chamber map and the scope boundary}
Let $H$ be the $D_4$ Cartan matrix. The chamber-ansatz map of
\cite[Def.~9]{KS} can be written
\[
 y_k=\prod_{\ell=1}^{12}x_\ell^{K_{k\ell}},\qquad
 K_{k\ell}=
 \begin{cases}
 1&\ell=k\text{ or }\ell=k^+\le12,\\
 H_{i_k,i_\ell}&k<\ell<k^+,\\
 0&\text{otherwise}.
 \end{cases}
\]
The upper-triangular matrix $K$ has diagonal one, hence determinant one.
A monomial row $a$ pulls back to $aK$. For every one of the $22$ rows in the
table, all entries of $aK$ lie in $\{-1,0,1\}$. In particular rows $21,22$
give, respectively,
\[
 (0,1,0,1,0,1,-1,0,0,0,0,0),\qquad
 (1,0,0,0,0,0,0,0,0,0,0,0).
\]
The monomial isomorphism therefore preserves the inequality presentation
but does not preserve its coordinatewise exponent bound.

The later $\varsigma$-coordinate converse in \cite[Conj.~3]{KS} is not
refuted: this example satisfies its exponent condition. Nor is its separate
coefficient criterion refuted, since all coefficients here are one and
all inequalities are essential. The computation contradicts the assertion
in the proof of \cite[Thm.~2, p.~1046]{KS} that the two exponent bounds are
equivalent under the chamber map. It does not contradict that theorem's
conclusion, and supplies no general replacement proof for its affected step.
Our direct facet proof does not use that assertion or that theorem.

\section{Prior computations and reproducibility}\label{sec:priority}
The primary-author dataset \cite[Case~267]{DATA} already records this word
and its complete $14+1+5+2$ string-coordinate support. The relevant Git commit
is dated 6 March 2020; a Git date is provenance metadata, not independent
proof of the date of public availability. The repository is cited by the
later computational work \cite{KM}, providing a public pre-PR citation trail.
We attribute the earlier word and support computations accordingly.
The published minuscule nonredundancy results \cite{OWR,KS} also cover the
three leaf-root summands; the central summand here is a single monomial.

The contribution of this note is to spell out the consequence for the
specified historical canonical-$W$ assertion, independently reconstruct the
canonical potential, and give complete directly checkable facet certificates.
A bounded literature audit found no earlier explicit statement of this
refutation, but cannot establish first priority, current openness, or
exhaustive coverage of inaccessible or unindexed sources. No claim is made
that the original converse remained open in October 2026.

The accompanying verification archive includes the manuscript source,
the complete exponent and certificate data, a Python standard-library
verifier that reconstructs the seed and optimized pullbacks over exact
integer Laurent arithmetic, and a separate SymPy reconstruction using
two reduced-word braid-graph searches. That reviewer inherited candidate
context; its freeze before a local file comparison does not establish
blindness. A further reviewer received only the word, Dynkin diagram and
primary references in a fresh context, froze its support before comparison,
and checked all 12,504 directed braid-graph edges and 22 separate facet
certificates. Its replay and original input-boundary receipt are included.
Source examples, mutation involutivity and deliberately corrupted inputs
provide implementation controls. The written formulas and displayed certificates, rather than
sampled numerical experiments, establish the result. Downloaded source
documents, third-party programs, and the original unlicensed dataset are
not redistributed.

\paragraph{AI use and review status.}
AI tools were used extensively to develop, solve, reconstruct, verify and
adversarially review the arguments, search the literature, and prepare this
note and its code. Independent internal AI reviews are not external human
peer review. This preprint is unrefereed and is not a formal
proof-assistant certificate.

\begin{thebibliography}{9}
\bibitem{OWR}
B.~Schumann (joint work with G.~Koshevoy),
Optimality of string cone inequalities and potential functions,
\emph{Oberwolfach Report} 20/2021, pp.~1123--1125.
\href{https://doi.org/10.4171/OWR/2021/20}{doi:10.4171/OWR/2021/20}.

\bibitem{KS}
G.~Koshevoy and B.~Schumann,
Redundancy in string cone inequalities and multiplicities in potential
functions on cluster varieties,
\emph{J.~Algebraic Combin.} \textbf{56} (2022), 1031--1053.
\href{https://doi.org/10.1007/s10801-022-01144-z}{doi:10.1007/s10801-022-01144-z}.

\bibitem{GKS}
V.~Genz, G.~Koshevoy and B.~Schumann,
Polyhedral parametrizations of canonical bases \& cluster duality,
\emph{Adv.~Math.} \textbf{369} (2020), 107178.
\href{https://doi.org/10.1016/j.aim.2020.107178}{doi:10.1016/j.aim.2020.107178}.

\bibitem{DATA}
D.~Mironov, \emph{crystal}, public computational repository,
\texttt{crystal\_D\_4.out.xz}, Case~267.
Pinned commit \texttt{0d6e564cf5e69d79c1074d5d1e3e72ac700904d6}
(commit dated 6 March 2020).
\url{https://github.com/mironovd/crystal}.

\bibitem{KM}
G.~Koshevoy and D.~Mironov,
$F$-polynomials and Newton polytopes,
\emph{EasyChair Preprint} 8975 (3 October 2022).
\url{https://easychair.org/publications/preprint/z55N}.
\end{thebibliography}
\end{document}
